% ================= Feature 002 (Lê Sĩ Nhật Quang) =================
\section{Feature 002: Add a new course -- kiểm thử hướng dữ liệu}
Mã nguồn nằm trong \texttt{F002/level1/}, \texttt{F002/level2/} và \texttt{F002/nonfunctional/} (Python 3 + Selenium 4 + unittest).

\subsection{Phân tích gộp 38 test case thành một script}
\begin{enumerate}[noitemsep]
  \item Ở bài tập lớn 2, 38 test case TC-002-001 $\dots$ 038 được ghi bằng Katalon Recorder (\texttt{F002-Add\_a\_new\_course.krecorder})
        cùng script dọn dữ liệu (\texttt{F002-Cleanup.krecorder}). Các test case này có \textbf{cùng một chuỗi thao tác} (mở form
        $\rightarrow$ nhập 7 trường $\rightarrow$ bấm nút $\rightarrow$ kiểm tra kết quả), chỉ khác nhau ở \emph{dữ liệu} và
        \emph{kết quả mong đợi}. Vì vậy cả 38 test case được gộp thành \textbf{một} kịch bản tự động.
  \item Toàn bộ dữ liệu và kết quả mong đợi của 38 test case được trích ra file \texttt{data/tc002\_add\_course.csv} (mỗi dòng là một test case).
  \item Script \texttt{test\_add\_course\_level1.py} đọc file dữ liệu và sinh một phương thức \texttt{unittest} cho mỗi dòng
        (\texttt{test\_TC\_002\_001}, \dots), nên mỗi test case vẫn có kết quả PASS/FAIL riêng.
  \item So với script Katalon, các locator động (\texttt{yui\_3\_18\_\dots}, \texttt{form\_autocomplete\_input-1794\dots},
        \texttt{course-info-container-7-4}) được thay bằng locator ổn định (\texttt{id=id\_fullname},
        \texttt{//input[contains(@id,'form\_autocomplete\_input')]}). Kết quả được kiểm tra ngay trên trang khóa học
        thay vì trang \textit{My courses}, nên script chạy lại được trên máy khác.
\end{enumerate}

\subsection{Cấu trúc file dữ liệu Level 1}
\begin{table}[H]
\centering\footnotesize
\begin{tabular}{|l|L{11.5cm}|}
\hline
\rowcolor{gray!15}\textbf{Cột} & \textbf{Ý nghĩa} \\ \hline
\texttt{tc\_id, technique, description} & Mã test case (trùng file \texttt{.xlsx} của BTL2), kỹ thuật (ECP, BVA, DT, UC), mô tả \\ \hline
\texttt{fullname, shortname, category, visibility} & Dữ liệu nhập cho các trường của form \\ \hline
\texttt{start, end, idnumber} & Ngày dạng \texttt{d/m/yyyy hh:mm}; \texttt{end} rỗng nghĩa là tắt \textit{Enable}; ID number \\ \hline
\texttt{action} & Nút bấm: \texttt{save\_display}, \texttt{save\_return}, \texttt{cancel} \\ \hline
\texttt{expected\_type} & \texttt{success} (tạo khóa học), \texttt{error} (báo lỗi tại trường), \texttt{message} (thông báo cấp trang, không tạo khóa học), \texttt{returned}, \texttt{cancelled} \\ \hline
\texttt{expected\_field} & Trường phải báo lỗi; nhiều trường ngăn cách bởi \texttt{|} (decision table R4, R5) \\ \hline
\texttt{expected\_result} & \textbf{Kết quả mong đợi}: thông báo lỗi, hoặc văn bản phải xuất hiện trên trang khóa học \\ \hline
\texttt{fix, fix\_expected\_result} & Cho use-case path ``lỗi $\rightarrow$ sửa $\rightarrow$ gửi lại'': dữ liệu sửa và kết quả mong đợi sau khi sửa \\ \hline
\texttt{entry} & Điểm vào form: rỗng = Site administration $>$ Add a new course; \texttt{category} = Manage courses and categories $>$ Create new course (chỉ điểm vào này có nút \textit{Save and return}, dùng cho TC-002-034) \\ \hline
\end{tabular}
\caption{Các cột của file dữ liệu}
\end{table}
\textbf{Macro trong dữ liệu:} \texttt{\{RUN\}} là mã duy nhất của lần chạy (tránh trùng short name khi chạy lại);
\texttt{\{fill:N:c\}} đệm chuỗi bằng ký tự \texttt{c} đến đúng \texttt{N} ký tự (dùng cho BVA 99, 100, 101, 253, 254, 255 ký tự);
\texttt{\{=fullname\}} nghĩa là kết quả mong đợi bằng giá trị của cột \texttt{fullname}.

Ví dụ ba dòng (rút gọn):
\begin{lstlisting}
tc_id,...,fullname,shortname,category,...,end,idnumber,action,expected_type,expected_field,expected_result
TC-002-003,...,,TC002-003-{RUN},Category 1,...,1/1/2027 08:00,,save_display,error,fullname,Missing full name
TC-002-017,...,{fill:254:e},TC002-017-{RUN},Category 1,...,1/1/2027 08:00,,save_display,success,,{=fullname}
TC-002-031,...,,,,...,1/1/2027 08:00,,save_display,error,fullname|shortname|category,Missing full name|Missing short name|You must supply a value here.
\end{lstlisting}

\subsection{Thiết kế script Level 1}
\begin{itemize}[noitemsep]
  \item \texttt{setUpClass}: mở Chrome, đăng nhập admin, \textbf{chuẩn bị dữ liệu}: tạo course \texttt{EXIST01}/\texttt{ID-EXIST}
        nếu chưa có (dùng cho TC-002-008, 013, 032, 037).
  \item \texttt{run\_case(row)}: mở \texttt{/course/edit.php?category=0}, điền 7 trường theo dòng dữ liệu, bấm nút ở cột \texttt{action}, rồi
        \textbf{verify}:
    \begin{itemize}[noitemsep]
      \item \texttt{success}: URL chuyển tới \texttt{/course/view.php} và \texttt{expected\_result} có trên trang.
      \item \texttt{error}: thẻ \texttt{id\_error\_<field>} chứa đúng thông báo trong \texttt{expected\_result}, form vẫn ở \texttt{edit.php}.
      \item \texttt{returned} / \texttt{cancelled}: tìm khóa học theo short name, phải có đúng một khóa học (Save and return) hoặc không có khóa học nào (Cancel).
    \end{itemize}
    Nếu có cột \texttt{fix}, script sửa dữ liệu, gửi lại và verify thành công.
  \item Ô nhập được gán giá trị trực tiếp giống lệnh \texttt{type} của Katalon (không bị \texttt{maxlength} cắt), để tái hiện đúng
        kết quả của TC-002-005, 007, 012 ở BTL2. Course category là ô autocomplete được JavaScript dựng trên thẻ
        \texttt{<select id="id\_category">}; script chọn trực tiếp trên thẻ select này (giống lệnh \texttt{select id=id\_category} của Katalon),
        với giá trị rỗng để tái hiện ``No selection'' (TC-002-009, 031) vì mặc định select đã chọn sẵn Category 1.
  \item Nút được cuộn vào giữa màn hình trước khi bấm vì footer cố định của Moodle che nút \textit{Save and display} khi form dài.
  \item \texttt{tearDown}: \textbf{dọn dữ liệu}, xóa mọi khóa học test vừa tạo; \texttt{tearDownClass} xóa course chuẩn bị và
        ghi \texttt{results\_level1.csv}.
\end{itemize}

\textbf{Đoạn code chính:}
\begin{lstlisting}[language=Python]
def run_case(self, row):
    self.page.open_form()
    self.page.fill(row)                 # 7 fields from the data row
    self.page.press(row["action"])      # save_display / save_return / cancel
    if row["expected_type"] == "success":
        self.check_success(row["expected_result"])
    elif row["expected_type"] == "error":
        self.check_errors(row["expected_field"], row["expected_result"])
        if row["fix"]:                  # use-case path: fix and resubmit
            self.page.apply_fix(row["fix"]); self.page.press(row["action"])
            self.check_success(row["fix_expected_result"])
    ...
for row in load_rows():                 # one unittest method per data row
    setattr(AddCourseDataDriven, "test_" + row["tc_id"].replace("-", "_"), _make_test(row))
\end{lstlisting}

\subsection{Level 2: đưa test item ra ngoài script}
Ở Level 2, script \texttt{test\_add\_course\_level2.py} là một \textbf{engine chung}, không chứa URL, id hay thông báo nào của Moodle.
Mọi \textit{test item} được đọc từ thư mục \texttt{config/}:
\begin{table}[H]
\centering\footnotesize
\begin{tabular}{|l|L{5.4cm}|L{6.2cm}|}
\hline
\rowcolor{gray!15}\textbf{File} & \textbf{Nội dung} & \textbf{Ví dụ} \\ \hline
\texttt{config/site.csv} & URL site, tài khoản, URL form/tìm kiếm/xóa, mẫu locator thông báo lỗi và ô ngày, dữ liệu chuẩn bị &
  \texttt{form\_url = \{base\_url\}/course/edit.php?category=0}; \texttt{error\_locator = id=id\_error\_\{field\}} \\ \hline
\texttt{config/locators.csv} & Tên, chiến lược (\texttt{id}, \texttt{xpath}, \texttt{css}, \dots) và giá trị của mọi phần tử &
  \texttt{save\_display, id, id\_saveanddisplay}; \texttt{category\_select, id, id\_category} \\ \hline
\texttt{config/fields.csv} & Cột dữ liệu nào điền vào phần tử nào và kiểu điều khiển (\texttt{text}, \texttt{select}, \texttt{hidden\_select}, \texttt{date}, \texttt{optional\_date}) &
  \texttt{end, optional\_date, enddate} \\ \hline
\texttt{data/tc002\_add\_course.csv} & 38 test case (giống Level 1), cột \texttt{action} là \emph{tên locator} của nút cần bấm & \texttt{action = save\_return} \\ \hline
\end{tabular}
\caption{Các file đầu vào của Level 2}
\end{table}
\textbf{Lợi ích:} khi Moodle đổi giao diện (id nút, URL) hoặc chạy trên site khác, chỉ cần sửa file cấu hình. Có thể chạy với
bộ dữ liệu hoặc cấu hình khác bằng tham số:
\begin{lstlisting}
cd F002/level2
python test_add_course_level2.py --data data/tc002_add_course.csv --config config
\end{lstlisting}
Engine có các thao tác chung: \texttt{set\_field(kind, locator, value)}, \texttt{press(locator)}, \texttt{error\_text(field)}
(dựng locator từ mẫu \texttt{error\_locator}), \texttt{find\_ids} và \texttt{delete} (dùng URL trong \texttt{site.csv}).
Quy trình verify, chuẩn bị và dọn dữ liệu giống Level 1.

\subsection{Cách chạy}
\begin{lstlisting}
pip install selenium
cd F002/level1
python test_add_course_level1.py                  # all 38 test cases
python test_add_course_level1.py -k TC_002_008    # one test case
\end{lstlisting}

\subsection{Kết quả chạy}
\begin{table}[H]\centering\small
\begin{tabular}{|L{4.6cm}|c|c|c|L{4.2cm}|}\hline
\rowcolor{gray!15}\textbf{Bộ test} & \textbf{Số test} & \textbf{Pass} & \textbf{Fail} & \textbf{Minh chứng} \\ \hline
Level 1 (data-driven) & 38 & 38 & 0 & \texttt{run\_level1.log}, \texttt{results\_level1.csv} \\ \hline
Level 2 (data + test items) & 38 & 38 & 0 & \texttt{run\_level2.log}, \texttt{results\_level2.csv} \\ \hline
NFR-1 Hiệu năng & 2 & 2 & 0 & \texttt{nf\_performance.log}, \texttt{perf\_results.csv} \\ \hline
NFR-2 Bảo mật & 3 & 3 & 0 & \texttt{nf\_security.log} \\ \hline
NFR-3 Tương thích & 4 & 4 & 0 & \texttt{nf\_compatibility.log} \\ \hline
\end{tabular}\caption{Kết quả chạy Feature 002 (Chrome, chế độ headless, Moodle demo 5.1)}\end{table}

{\footnotesize\begin{longtable}{|l|c|L{1.4cm}|L{6.0cm}|c|c|}\hline
\rowcolor{gray!15}\textbf{ID} & \textbf{KT} & \textbf{Loại} & \textbf{Kết quả mong đợi (từ file dữ liệu)} & \textbf{L1} & \textbf{L2} \\ \hline\endhead
TC-002-001 & ECP & success & = fullname & PASS & PASS \\ \hline
TC-002-002 & ECP & success & = fullname & PASS & PASS \\ \hline
TC-002-003 & ECP & error & Missing full name & PASS & PASS \\ \hline
TC-002-004 & ECP & success & Course: \textbar{} & PASS & PASS \\ \hline
TC-002-005 & ECP/\allowbreak{}BVA & success & = fullname & PASS & PASS \\ \hline
TC-002-006 & ECP & error & Missing short name & PASS & PASS \\ \hline
TC-002-007 & ECP/\allowbreak{}BVA & success & = fullname & PASS & PASS \\ \hline
TC-002-008 & ECP & error & Short name is already used for another course & PASS & PASS \\ \hline
TC-002-009 & ECP & error & You must supply a value here. & PASS & PASS \\ \hline
TC-002-010 & ECP/\allowbreak{}BVA & success & = fullname & PASS & PASS \\ \hline
TC-002-011 & ECP & error & The course end date must be after the start date. & PASS & PASS \\ \hline
TC-002-012 & ECP/\allowbreak{}BVA & message & Error writing to database & PASS & PASS \\ \hline
TC-002-013 & ECP & error & ID number is already used for another course & PASS & PASS \\ \hline
TC-002-014 & BVA & success & Course: A \textbar{} & PASS & PASS \\ \hline
TC-002-015 & BVA & success & Course: BB \textbar{} & PASS & PASS \\ \hline
TC-002-016 & BVA & success & = fullname & PASS & PASS \\ \hline
TC-002-017 & BVA & success & = fullname & PASS & PASS \\ \hline
TC-002-018 & BVA & success & = fullname & PASS & PASS \\ \hline
TC-002-019 & BVA & success & = fullname & PASS & PASS \\ \hline
TC-002-020 & BVA & success & = fullname & PASS & PASS \\ \hline
TC-002-021 & BVA & success & = fullname & PASS & PASS \\ \hline
TC-002-022 & BVA & success & = fullname & PASS & PASS \\ \hline
TC-002-023 & BVA & success & = fullname & PASS & PASS \\ \hline
TC-002-024 & BVA & success & = fullname & PASS & PASS \\ \hline
TC-002-025 & BVA & success & = fullname & PASS & PASS \\ \hline
TC-002-026 & BVA & error & The course end date must be after the start date. & PASS & PASS \\ \hline
TC-002-027 & BVA & success & = fullname & PASS & PASS \\ \hline
TC-002-028 & DT & success & = fullname & PASS & PASS \\ \hline
TC-002-029 & DT & success & = fullname & PASS & PASS \\ \hline
TC-002-030 & DT & error & The course end date must be after the start date. & PASS & PASS \\ \hline
TC-002-031 & DT & error & Missing full name\textbar{}Missing short name\textbar{}You must supply a value here. & PASS & PASS \\ \hline
TC-002-032 & DT & error & Short name is already used for another course\textbar{}The course end date must be after \dots{} & PASS & PASS \\ \hline
TC-002-033 & UC & success & = fullname & PASS & PASS \\ \hline
TC-002-034 & UC & returned & = fullname & PASS & PASS \\ \hline
TC-002-035 & UC & cancelled & = fullname & PASS & PASS \\ \hline
TC-002-036 & UC & error & Missing full name\textbar{}Missing short name $\rightarrow$ sửa $\rightarrow$ TC002 UC Retry & PASS & PASS \\ \hline
TC-002-037 & UC & error & Short name is already used for another course $\rightarrow$ sửa $\rightarrow$ TC002 UC Dup & PASS & PASS \\ \hline
TC-002-038 & UC & error & The course end date must be after the start date. $\rightarrow$ sửa $\rightarrow$ TC002 UC Date & PASS & PASS \\ \hline
\end{longtable}}

\subsection{Phát hiện khi tự động hóa lại}\label{sec:f002-findings}
Khi chạy lại bằng Python, hai kết quả của BTL2 (ghi từ Katalon) được phát hiện là \textbf{không đúng với hành vi thật} của Moodle
và đã được sửa trong file dữ liệu; cần cập nhật lại file \texttt{.xlsx} của BTL2 tương ứng:
\begin{table}[H]
\centering\footnotesize
\begin{tabular}{|l|L{3.4cm}|L{4.0cm}|L{5.2cm}|}
\hline
\rowcolor{gray!15}\textbf{Test case} & \textbf{Đầu vào} & \textbf{Ghi nhận ở BTL2} & \textbf{Hành vi thật (BTL3)} \\ \hline
TC-002-010 & End date = start date (1/1/2026 08:00) & Báo lỗi ``The course end date must be after the start date.'' &
  \textbf{Khóa học được tạo}: Moodle chỉ từ chối khi end $<$ start. Script Katalon cũ chỉ đặt ngày/tháng của end date, giờ vẫn
  mặc định 00:00 nên thực chất end $<$ start. Biên đúng: $-1$ phút lỗi (TC-002-026); $0$ và $+1$ phút hợp lệ. \\ \hline
TC-002-012 & ID number 101 ký tự & Form được lưu & \textbf{Lỗi hệ thống} ``Error writing to database -- Data too long for column idnumber'':
  form không kiểm tra độ dài phía server, cơ sở dữ liệu từ chối. Đây là \textbf{lỗi (defect)} nên báo cáo, mức độ trung bình. \\ \hline
\end{tabular}
\caption{Kết quả BTL2 được hiệu chỉnh nhờ tự động hóa lại}
\end{table}
Các hành vi khác đặc tả đã ghi nhận ở BTL2 được xác nhận lại: full name chỉ có khoảng trắng (TC-002-004), full name 255 ký tự (005),
short name 101 ký tự (007) đều được chấp nhận khi gán giá trị trực tiếp (bỏ qua \texttt{maxlength} của trình duyệt).
